% Self-contained mathematical section; included by the cumulative main.tex.
% Requires amsmath, amssymb, amsthm and the standard theorem environments.
\section{The exact axisymmetric coordinate mechanism and its viscous extension}
\label{sec:exact-euler-coordinates}

This section gives the complete coordinate and force reconstruction underlying
the Euler system on pages 3--4, equations (2.1)--(2.6), of
\emph{Blowup for the Euler equations with smooth forcing}, the released
manuscript examined in this repository.  The coordinate identities are proved
here independently.  The compact force inverse is constructed and proved
explicitly; no external inverse-divergence lemma is used.  The same calculation
also gives every additional term for classical viscosity \(\nu>0\).
It proves an exact correspondence of smooth equations.  It does not verify
the manuscript's subsequent infinite construction or its terminal estimates.

\subsection{Domains, signs, and the forward and inverse maps}

Fix the original physical parameters
\begin{gather*}
 r_0>0,\qquad z_0\in\mathbb R,\qquad 0<R<r_0/2,
 \qquad r=(x_1^2+x_2^2)^{1/2},\\
 \mathcal T_R=\{(x_1,x_2,z):(z-z_0)^2+(r-r_0)^2<R^2\}.
\end{gather*}
No dilation or translation of these parameters is performed.  On \(r>0\),
write
\[
 e_r=(\cos\phi,\sin\phi,0),\qquad
 e_\phi=(-\sin\phi,\cos\phi,0),\qquad e_z=(0,0,1).
\]
Thus \((e_r,e_\phi,e_z)\) is positively oriented.  An axisymmetric vector
field has the form
\(u=u^r(r,z,t)e_r+u^\phi(r,z,t)e_\phi+u^z(r,z,t)e_z\),
with components independent of \(\phi\).

The ordered meridional coordinates in this section are \((z,r)\).
Define the diffeomorphism and its inverse
\begin{align}
 \Phi:\mathbb R\times(0,\infty)&\longrightarrow
              H:=\mathbb R\times(0,\infty),
 &\Phi(z,r)&=(y_1,y_2)=(z,r^2/2),\label{eq:ec-Phi}\\
 \Phi^{-1}(y_1,y_2)&=(y_1,\sqrt{2y_2}).\nonumber
\end{align}
Their Jacobian determinants, in these orders, are \(r\) and \(1/r\).
In particular
\begin{gather*}
 dy_1\wedge dy_2=r\,dz\wedge dr,\qquad dy=r\,dz\,dr,\\
 \partial_z=\partial_{y_1},\qquad
 \partial_r=r\partial_{y_2},\qquad
 \partial_{y_2}=r^{-1}\partial_r.
\end{gather*}
For a meridional set \(K\subset H\), its physical lift is
\[
 \widetilde K=\{(r\cos\phi,r\sin\phi,z):
                 (z,r^2/2)\in K,\ \phi\in\mathbb R/(2\pi\mathbb Z)\}.
\]
Let \(D_0\Subset D\) be concentric open balls in \(H\), centered at
\(y_*=(z_0,r_0^2/2)\), whose closures lie in the image under \(\Phi\) of
the physical meridional disk defining \(\mathcal T_R\).  Such balls exist:
that image is open and contains \(y_*\), so it contains a sufficiently small
closed ball and a smaller concentric closed ball.  The notation
\(D_0\Subset D\) means that \(\overline{D_0}\) is a compact subset of \(D\).

All fields below are smooth on a time interval \(I\); at a closed endpoint,
smoothness means compatible one-sided derivatives.  Their spatial supports
are contained in the indicated fixed compact subset of the stated ball,
uniformly in \(t\in I\).  Fields in volume coordinates supported in \(H\)
are extended by zero to \(\mathbb R^2\).

Set
\begin{equation}
 \begin{gathered}
 J(a_1,a_2)=(-a_2,a_1),\qquad
 A(y)=\begin{pmatrix}(2y_2)^{-1}&0\\0&1\end{pmatrix},\\
 L=\partial_{y_2}^2+(2y_2)^{-1}\partial_{y_1}^2,
 \qquad W=(2y_2)^{-2}.
 \end{gathered}
 \label{eq:ec-coefficients}
\end{equation}
The equality \(L=\operatorname{div}_y(A\nabla_y)\) follows by expansion:
the first diagonal coefficient depends on \(y_2\), whereas the derivative
in its divergence term is \(\partial_{y_1}\); the other diagonal coefficient
is identically one.

\begin{proposition}[Velocity correspondence with a fixed streamfunction gauge]
\label{prop:ec-velocity}
The linear map from pairs \((\Gamma,\Psi)\) of smooth functions supported
in \(D_0\) to physical fields, defined by
\begin{equation}
 u^z=-\Psi_{y_2},\qquad
 u^r=\frac{\Psi_{y_1}}{r},\qquad
 u^\phi=\frac{\Gamma}{r},\qquad r=\sqrt{2y_2},
 \label{eq:ec-physical-velocity}
\end{equation}
is a bijection onto the smooth, divergence-free, axisymmetric vector fields
supported in \(\widetilde{D_0}\).  The inverse recovers
\begin{equation}
 \begin{gathered}
 \Gamma(y,t)=r u^\phi(r,z,t),\qquad
 v(y,t)=(u^z(r,z,t),r u^r(r,z,t)),\\
 \Psi(y,t)=\int_{-\infty}^{y_1}v_2(s,y_2,t)\,ds.
 \end{gathered}
 \label{eq:ec-inverse}
\end{equation}
In particular \(v=J\nabla_y\Psi\).
\end{proposition}

\begin{proof}
For an axisymmetric field,
\begin{equation}
 \operatorname{div}_x u
 =\partial_z u^z+r^{-1}\partial_r(r u^r)
 =\partial_{y_1}v_1+\partial_{y_2}v_2.
 \label{eq:ec-divergence}
\end{equation}
One direct derivation of the first identity uses
\(\partial_{x_1}r=\cos\phi\),
\(\partial_{x_2}r=\sin\phi\),
\(\partial_{x_1}\phi=-\sin\phi/r\), and
\(\partial_{x_2}\phi=\cos\phi/r\), substituted into the Cartesian
divergence.  The terms containing the azimuthal component cancel; those
containing the radial component give \(\partial_r u^r+u^r/r\).
For \eqref{eq:ec-physical-velocity}, \(v=(-\Psi_{y_2},\Psi_{y_1})\), so
the right side of \eqref{eq:ec-divergence} is zero by equality of mixed
derivatives.  The factors \(1/r\), the coordinate change, and the
cylindrical basis are smooth on a neighborhood of the fixed support.
The field is identically zero on a fixed neighborhood of the axis.
Consequently its extension to all of \(\mathbb R^3\) is smooth, and its
support is as stated.

Conversely, let \(u\) be a field in the specified target space.  The
functions \(\Gamma\) and \(v\) in \eqref{eq:ec-inverse} are smooth and
compactly supported in \(D_0\).  Equation \eqref{eq:ec-divergence} gives
\(\partial_{y_1}v_1+\partial_{y_2}v_2=0\).  The integral in
\eqref{eq:ec-inverse} is smooth and satisfies
\[
 \Psi_{y_1}=v_2,\qquad
 \Psi_{y_2}=\int_{-\infty}^{y_1}\partial_{y_2}v_2(s,y_2,t)\,ds
 =-\int_{-\infty}^{y_1}\partial_s v_1(s,y_2,t)\,ds=-v_1.
\]
Here the lower boundary term vanishes by compact support.
It remains to verify compact support of \(\Psi\), which cannot simply
be inferred from compact support of its gradient.  Define
\(m(y_2,t)=\int_{\mathbb R}v_2(s,y_2,t)\,ds\).  Its derivative is
\[
 \partial_{y_2}m=-\int_{\mathbb R}\partial_s v_1(s,y_2,t)\,ds=0.
\]
It vanishes for \(y_2\) outside the projection of the support, hence
\(m=0\) everywhere.  Each horizontal slice of the ball \(D_0\) is
an interval.  To the left of that interval the integral defining \(\Psi\)
is zero; to its right it equals \(m=0\).  Thus \(\Psi=0\) outside
\(D_0\).  A smaller concentric ball still contains the fixed compact
support of \(v\), and the same argument places the support of \(\Psi\)
in that smaller ball.  Therefore \(\Psi\) has the required compact support.
Any two streamfunctions giving this \(v\) have zero gradient, so their
difference is spatially constant.  Compact support makes that constant
zero.  The formulas recover all three components of \(u\), proving the
bijection, including its gauge and its time dependence.
\end{proof}

For every smooth axisymmetric scalar \(a\), the material derivative is
preserved exactly:
\begin{equation}
 D a:=(\partial_t+u^r\partial_r+u^z\partial_z)a
   =(\partial_t+v\cdot\nabla_y)a=:D_v a.
 \label{eq:ec-material}
\end{equation}
This identity applies to scalar components; the derivative of a vector
also differentiates the cylindrical basis, as computed below.
The measure identity also gives the exact energy formula
\begin{align}
 \int_{\mathbb R^3}|u(x,t)|^2\,dx
 &=2\pi\int_H\left[
       \Psi_{y_2}^2+\frac{\Psi_{y_1}^2+\Gamma^2}{2y_2}\right]dy
 \label{eq:ec-energy}\\
 &=2\pi\int_H\left[v_1^2+\frac{v_2^2+\Gamma^2}{2y_2}\right]dy.
 \nonumber
\end{align}
Indeed the azimuthal integral of a scalar independent of \(\phi\) is
\(2\pi\), and \(r\,dr\,dz=dy\); insertion of
\eqref{eq:ec-physical-velocity} proves both equalities.

\subsection{Full vorticity and the exact physical gradient correspondence}

The chosen orientation gives
\begin{equation}
 \omega^r=-\partial_z u^\phi,\qquad
 \omega^\phi=\partial_z u^r-\partial_r u^z,\qquad
 \omega^z=r^{-1}\partial_r(r u^\phi),
 \qquad \omega=\nabla_x\times u.
 \label{eq:ec-curl-cylindrical}
\end{equation}
These identities can be checked without a coordinate convention left
implicit.  At \(\phi=0\),
\(\partial_{x_1}=\partial_r\) and
\(\partial_{x_2}=r^{-1}\partial_\phi\).
The basis derivatives are
\(\partial_\phi e_r=e_\phi\) and
\(\partial_\phi e_\phi=-e_r\).
Substitution into the Cartesian curl therefore gives, in order,
\(-\partial_z u^\phi\),
\(\partial_z u^r-\partial_r u^z\), and
\(\partial_r u^\phi+u^\phi/r\).
Rotational covariance gives the same components at every \(\phi\).

Define the reduced azimuthal vorticity and the squared circulation by
\[
 \xi=\frac{\omega^\phi}{r},\qquad \Theta=\Gamma^2.
\]
The velocity map gives every component, with its original radial factor:
\begin{equation}
 \omega^r=-\frac{\Gamma_{y_1}}r,\qquad
 \omega^\phi=rL\Psi=r\xi,\qquad
 \omega^z=\Gamma_{y_2},\qquad \xi=L\Psi.
 \label{eq:ec-full-vorticity}
\end{equation}
For the middle component, direct differentiation gives
\[
 r^{-1}(\partial_z u^r-\partial_r u^z)
 =r^{-1}\left(\frac{\Psi_{y_1y_1}}r
                          +r\Psi_{y_2y_2}\right)
 =(2y_2)^{-1}\Psi_{y_1y_1}+\Psi_{y_2y_2}.
\]
The first and third components follow by differentiating \(\Gamma/r\)
in \eqref{eq:ec-curl-cylindrical}.  Consequently the full physical
gradient, rather than just its coordinate representation, obeys
\begin{equation}
 \nabla_x\Gamma=r\Gamma_{y_2}e_r+\Gamma_{y_1}e_z,
 \qquad
 |\omega|^2=\frac{|\nabla_x\Gamma|^2}{r^2}+r^2\xi^2.
 \label{eq:ec-gradient-vorticity}
\end{equation}
The two displayed vectors contributing to \(\omega\) are orthogonal,
which proves the sum of squares.
For fields supported in \(\mathcal T_R\), put
\(r_-=r_0-R>0\) and \(r_+=r_0+R\).  Pointwise on their support,
\[
 |\omega|\geq r_+^{-1}|\nabla_x\Gamma|,
 \qquad
 |\omega|\leq r_-^{-1}|\nabla_x\Gamma|+r_+|\xi|.
\]
Taking suprema preserves these inequalities.  Thus growth of the
physical circulation gradient forces full vorticity growth by an
explicit bounded geometric factor.  This statement alone gives no
growth rate or time-integrability conclusion.

\subsection{Derivation of the reduced Euler and Navier--Stokes equations}

Keep \(\nu\geq0\) as an explicit constant and use the convention
\begin{equation}
 \partial_tu+(u\cdot\nabla_x)u+\nabla_xp
      =\nu\Delta_xu+f,\qquad \operatorname{div}_xu=0.
 \label{eq:ec-NS}
\end{equation}
The value \(\nu=0\) is Euler.  The force and pressure are axisymmetric
in this calculation.  Define the scalar cylindrical Laplacian
\(\Delta_0=\partial_z^2+\partial_r^2+r^{-1}\partial_r\).
The basis derivatives above, together with
\(u\cdot\nabla=u^r\partial_r+(u^\phi/r)\partial_\phi+u^z\partial_z\),
give the acceleration components
\[
 D u^r-\frac{(u^\phi)^2}{r},\qquad
 D u^\phi+\frac{u^r u^\phi}{r},\qquad D u^z.
\]
For completeness the scalar Laplacian follows by applying the divergence
formula to the scalar gradient:
\(\Delta a=r^{-1}\partial_r(r\partial_ra)
+r^{-2}\partial_\phi^2a+\partial_z^2a\).
Apply this Cartesian-componentwise to a vector.  Since
\(\partial_\phi^2e_r=-e_r\),
\(\partial_\phi^2e_\phi=-e_\phi\), and
\(\partial_\phi e_z=0\), its axisymmetric components are
\[
 (\Delta_xu)^r=(\Delta_0-r^{-2})u^r,\qquad
 (\Delta_xu)^\phi=(\Delta_0-r^{-2})u^\phi,
 \qquad (\Delta_xu)^z=\Delta_0u^z.
\]
Thus \eqref{eq:ec-NS} is precisely
\begin{align}
 D u^r-\frac{(u^\phi)^2}{r}+p_r
       &=\nu(\Delta_0-r^{-2})u^r+f^r,\label{eq:ec-radial}\\
 D u^\phi+\frac{u^r u^\phi}{r}
       &=\nu(\Delta_0-r^{-2})u^\phi+f^\phi,\label{eq:ec-swirl}\\
 D u^z+p_z&=\nu\Delta_0u^z+f^z.\label{eq:ec-axial}
\end{align}

Multiplying \eqref{eq:ec-swirl} by \(r\) and using \(Dr=u^r\) gives
\begin{equation}
 D\Gamma=\nu\mathcal D_-\Gamma+rf^\phi,
 \qquad \mathcal D_-:=\partial_z^2+\partial_r^2-r^{-1}\partial_r.
 \label{eq:ec-Gamma-physical}
\end{equation}
To verify the diffusion term explicitly,
\[
 \partial_r(\Gamma/r)=\Gamma_r/r-\Gamma/r^2,\qquad
 \partial_r^2(\Gamma/r)=\Gamma_{rr}/r-2\Gamma_r/r^2+2\Gamma/r^3,
\]
so \(r(\Delta_0-r^{-2})(\Gamma/r)
=\Gamma_{zz}+\Gamma_{rr}-\Gamma_r/r\).

Next set \(q=\partial_z u^r-\partial_r u^z=r\xi\).
For arbitrary smooth \(u^r,u^z\), expansion of the material derivative
gives
\begin{align*}
 \partial_z(Du^r)-\partial_r(Du^z)
 &=Dq+(\partial_zu^r)(\partial_ru^r)
       +(\partial_zu^z)(\partial_zu^r)\\
 &\hspace{9mm}-(\partial_ru^r)(\partial_ru^z)
       -(\partial_ru^z)(\partial_zu^z)\\
 &=Dq+(\partial_ru^r+\partial_zu^z)q
 =Dq-\frac{u^r}{r}q=rD\xi.
\end{align*}
The penultimate equality uses exactly the divergence equation, and the
last uses \(D(r\xi)=u^r\xi+rD\xi\).
Apply \(\partial_z\) to \eqref{eq:ec-radial} and subtract
\(\partial_r\) of \eqref{eq:ec-axial}.  The pressure cancels by equality
of mixed derivatives.  The centrifugal term becomes
\[
 -\partial_z((u^\phi)^2/r)=-r^{-3}\partial_z(\Gamma^2).
\]
The viscous term is computed from the explicit commutator
\[
 \partial_r(\Delta_0 a)=\Delta_0(\partial_ra)-r^{-2}\partial_ra,
 \qquad \partial_z\Delta_0a=\Delta_0\partial_za.
\]
These give
\begin{align*}
 \partial_z[(\Delta_0-r^{-2})u^r]-\partial_r\Delta_0u^z
  &=(\Delta_0-r^{-2})q\\
  &=(\Delta_0-r^{-2})(r\xi)
   =r\left(\xi_{zz}+\xi_{rr}+\frac3r\xi_r\right).
\end{align*}
In the last line, \((r\xi)_r=\xi+r\xi_r\) and
\((r\xi)_{rr}=2\xi_r+r\xi_{rr}\), so all zero-order terms cancel
exactly.  Division by the original radius now proves
\begin{equation}
 D\xi=r^{-4}\partial_z(\Gamma^2)
       +\nu\mathcal D_+\xi
       +\frac{\partial_zf^r-\partial_rf^z}{r},
 \qquad \mathcal D_+:=\partial_z^2+\partial_r^2+3r^{-1}\partial_r.
 \label{eq:ec-xi-physical}
\end{equation}

For any scalar \(a(y_1,y_2,t)\),
\[
 \partial_r a=r a_{y_2},\qquad
 \partial_r^2 a=a_{y_2}+r^2a_{y_2y_2},\qquad
 \partial_z^2a=a_{y_1y_1}.
\]
Therefore the two distinct diffusion operators in volume coordinates are
\begin{equation}
 Q:=\partial_{y_1}^2+2y_2\partial_{y_2}^2,
 \qquad M:=\partial_{y_1}^2+2y_2\partial_{y_2}^2+4\partial_{y_2},
 \qquad \mathcal D_-=Q,\quad\mathcal D_+=M.
 \label{eq:ec-volume-diffusion}
\end{equation}
In particular \(Q=2y_2L\).  This is equality of operators applied to a
function, with multiplication by \(2y_2\) performed \emph{after} applying
\(L\); it does not permit commuting derivatives through that factor.

For the meridional force use the covector components
\begin{equation}
 b=(b_1,b_2)=(f^z,f^r/r),\qquad
 \operatorname{curl}_y b:=\partial_{y_1}b_2-\partial_{y_2}b_1.
 \label{eq:ec-force-covector}
\end{equation}
They differ from the pushforward velocity components because
\[
 f^z\,dz+f^r\,dr=b_1\,dy_1+b_2\,dy_2,
 \qquad
 \operatorname{curl}_yb
       =\frac{\partial_zf^r-\partial_rf^z}{r}.
\]
The differential of this one-form is
\((\operatorname{curl}_yb)dy_1\wedge dy_2\), which agrees with
\((\partial_zf^r-\partial_rf^z)dz\wedge dr\) by the Jacobian.
This proves the exact map for the force, including its orientation.

The complete reduced system is consequently
\begin{equation}
 \boxed{\begin{aligned}
    v&=J\nabla_y\Psi,\qquad \xi=L\Psi,\\
    D_v\Gamma&=\nu Q\Gamma+rf^\phi,\\
    D_v\xi&=W\partial_{y_1}(\Gamma^2)+\nu M\xi
                          +\operatorname{curl}_y(f^z,f^r/r).
 \end{aligned}}
 \label{eq:ec-reduced-NS}
\end{equation}
The equation for the original square \(\Theta=\Gamma^2\), with its
gradient-square term retained, is
\begin{equation}
 D_v\Theta=\nu Q\Theta
 -2\nu\bigl((\Gamma_{y_1})^2+2y_2(\Gamma_{y_2})^2\bigr)
 +2\Gamma r f^\phi.
 \label{eq:ec-Theta}
\end{equation}
Indeed the product rule yields
\(Q(\Gamma^2)=2\Gamma Q\Gamma
+2(\Gamma_{y_1})^2+4y_2(\Gamma_{y_2})^2\), and multiplication of the
circulation equation by \(2\Gamma\) gives \eqref{eq:ec-Theta}.
When \(\nu=0\), \eqref{eq:ec-reduced-NS} and
\eqref{eq:ec-Theta} are exactly the Euler formulas (2.6) in the source.

The momentum equations can also be displayed directly in volume
coordinates, preserving the geometric quadratic terms:
\begin{align}
 D_vv_1+p_{y_1}&=f^z+\nu(Q+2\partial_{y_2})v_1,\label{eq:ec-v1}\\
 D_vv_2-\frac{v_2^2+\Gamma^2}{2y_2}+2y_2p_{y_2}
      &=rf^r+\nu Qv_2.\label{eq:ec-v2}
\end{align}
The first follows from \(v_1=u^z\) and
\(\Delta_0=Q+2\partial_{y_2}\).  For the second,
\(D(r u^r)=(u^r)^2+rDu^r\); multiply
\eqref{eq:ec-radial} by \(r\), use \(r p_r=r^2p_{y_2}\),
and apply the already proved identity
\(r(\Delta_0-r^{-2})(v_2/r)=Qv_2\).

\subsection{An explicit compact curl inverse, with proof and bounds}

We construct the force operator used below.  Choose once and for all a
smooth nonnegative function \(\eta\) supported in a compact subset of
\(D\), with \(\int_{\mathbb R^2}\eta=1\).  Existence is explicit:
choose a closed ball inside \(D\), put
\(\eta_0(x)=\exp[-1/(1-|x-a|^2/\rho^2)]\) for \(|x-a|<\rho\)
and zero otherwise, and divide it by its positive finite integral.
Every derivative at the boundary vanishes because a polynomial in
\((1-|x-a|^2/\rho^2)^{-1}\) times the exponential tends to zero.
This proves the required smoothness of the bump.

For a smooth \(g\) supported in \(D_0\), define
\begin{equation}
 (Bg)(x)=\int_{\mathbb R^2}g(y)(x-y)
       \left(\int_1^\infty\eta(y+t(x-y))t\,dt\right)dy,
 \qquad \mathcal Rg=JBg.
 \label{eq:ec-B}
\end{equation}
The formula is interpreted by its absolutely locally integrable vector
kernel, as justified next; equivalently it is given by the convergent
smooth formulas in the proof.

\begin{lemma}[Fixed compact inverse]
\label{lem:ec-inverse-div}
The operator \(B\) is linear and smooth on smooth compactly supported
inputs, and
\begin{equation}
 \operatorname{div}Bg=g-\eta\int_{\mathbb R^2}g,
 \qquad
 \operatorname{curl}_y\mathcal Rg=g-\eta\int_{\mathbb R^2}g.
 \label{eq:ec-B-div}
\end{equation}
Its support is contained in the convex hull of
\(\operatorname{supp}g\cup\operatorname{supp}\eta\), hence compactly
contained in \(D\).  For every nonnegative integer \(k\) there is a
finite constant \(C_k\), depending only on the fixed balls, the bump,
and \(k\), such that
\begin{equation}
 \|Bg\|_{C^k(\mathbb R^2)}+
 \|\mathcal Rg\|_{C^k(\mathbb R^2)}\leq C_k\|g\|_{C^k(\mathbb R^2)}.
 \label{eq:ec-B-bound}
\end{equation}
The same operator commutes with every time derivative.
\end{lemma}

\begin{proof}
Put \(s=t-1\) and \(w=x-y\) in \eqref{eq:ec-B}.  Split the resulting
\(s\)-integral at one.  The portion from zero to one is
\begin{equation}
 B_{\rm near}g(x)=\int_0^1(1+s)
              \int_{\mathbb R^2}w g(x-w)\eta(x+sw)\,dw\,ds.
 \label{eq:ec-B-near}
\end{equation}
In the portion from one to infinity put \(a=sw\).  It is
\begin{equation}
 B_{\rm far}g(x)=\int_1^\infty\frac{1+s}{s^3}
              \int_{\mathbb R^2}a g(x-a/s)\eta(x+a)\,da\,ds.
 \label{eq:ec-B-far}
\end{equation}
For \(x\) in any fixed bounded set, nonzero integrands in
\eqref{eq:ec-B-near} have \(w\) in a fixed bounded set because
\(x-w\in\operatorname{supp}g\subset\overline{D_0}\).
The integration domains in \(w,s\) then have finite measure and all
\(x\)-derivatives are bounded by fixed finite linear combinations of
derivatives of \(g\) and \(\eta\).  In \eqref{eq:ec-B-far},
\(a\) lies in a fixed bounded set because
\(x+a\in\operatorname{supp}\eta\).  The scalar weight is integrable:
\(\int_1^\infty(1+s)s^{-3}ds=3/2\).  Differentiating with respect to
\(x\) differentiates only \(g(x-a/s)\) and \(\eta(x+a)\), with
coefficients one.  Dominated convergence therefore permits every
derivative, proves smoothness, and bounds derivatives through order
\(k\) by a constant times \(\|g\|_{C^k}\).  These arguments also
prove absolute convergence of the original vector integral locally in
\(x\), by the indicated changes of variables.

If an integrand in \eqref{eq:ec-B} is nonzero, then
\(y\in\operatorname{supp}g\) and
\(a=y+t(x-y)\in\operatorname{supp}\eta\) for some \(t\geq1\).
Consequently \(x=(1-1/t)y+(1/t)a\) lies on the segment joining those
two compact supports.  Outside their closed convex hull the integral
vanishes.  Both supports lie in a smaller concentric closed ball inside
the convex ball \(D\), so their convex hull is a compact subset of
\(D\).  Apply the preceding derivative bounds on one fixed bounded
neighborhood of \(\overline D\); outside that neighborhood the function
and its derivatives vanish.  This proves the global estimate
\eqref{eq:ec-B-bound}, also for \(JB\) because \(J\) preserves
Euclidean norm.

To prove the divergence identity, let \(\varphi\) be a smooth compactly
supported test function.  In the pairing with \(\nabla\varphi\), change
variables \(a=y+t(x-y)\) in the \(x\)-integral.  Then
\(dx=t^{-2}da\) and \(x-y=(a-y)/t\).  Thus
\begin{align*}
 \int Bg(x)\cdot\nabla\varphi(x)\,dx
  &=\int g(y)\int\eta(a)\int_1^\infty
       \frac{a-y}{t^2}\cdot
        \nabla\varphi\left(y+\frac{a-y}{t}\right)dt\,da\,dy\\
  &=\int g(y)\int\eta(a)\,[\varphi(a)-\varphi(y)]\,da\,dy.
\end{align*}
All integrations here are absolutely convergent: \(a,y\) range over
fixed compact sets and the remaining bound is a constant times
\(t^{-2}\), whose integral is finite.  The second equality uses
\[
 \frac{d}{dt}\varphi\left(y+\frac{a-y}{t}\right)
 =-\frac{a-y}{t^2}\cdot\nabla\varphi\left(y+\frac{a-y}{t}\right),
\]
and the limit at infinity is \(\varphi(y)\).
Since \(\int\eta=1\), the definition of distributional divergence
gives \(\operatorname{div}Bg=g-\eta\int g\).
Both sides are smooth, so equality also holds pointwise.  Finally
\(\operatorname{curl}_y Jb
=\partial_{y_1}b_1+\partial_{y_2}b_2=\operatorname{div}b\),
proving the curl identity with the stated sign.  The operator has no
time-dependent coefficients, and the same dominated-convergence
argument commutes any time derivative through it.  This proves all
claims of the lemma.
\end{proof}

\subsection{Force and pressure recovery in both directions}

For an arbitrary pair \((\Gamma,\Psi)\) as above define its exact
reduced residuals at viscosity \(\nu\) by
\begin{equation}
 F_\nu^\Gamma=D_v\Gamma-\nu Q\Gamma,\qquad
 E_\nu^\xi=D_v(L\Psi)-W\partial_{y_1}(\Gamma^2)-\nu M(L\Psi).
 \label{eq:ec-residuals}
\end{equation}
No evolution equation has been assumed in this definition.

\begin{proposition}[Complete compact reconstruction]
\label{prop:ec-reconstruction}
For every such smooth pair and every fixed \(\nu\geq0\), the formulas
\begin{equation}
 b=\mathcal R E_\nu^\xi,\qquad
 f^z=b_1,\qquad f^r=rb_2,\qquad f^\phi=F_\nu^\Gamma/r
 \label{eq:ec-force-recovery}
\end{equation}
give a smooth axisymmetric force supported in \(\widetilde D\).
The velocity \eqref{eq:ec-physical-velocity} then solves
\eqref{eq:ec-NS} with a smooth axisymmetric pressure on
\(\mathbb R^3\times I\).  All statements include mixed space-time
derivatives.  Conversely, every smooth axisymmetric solution of
\eqref{eq:ec-NS} in this compact velocity class satisfies
\eqref{eq:ec-reduced-NS} with its original force.  Hence the reduced
equations plus the proved pressure reconstruction are an exact
correspondence of the stated systems.
\end{proposition}

\begin{proof}
First we verify the zero integral required by the compact inverse.
All integrals are over \(\mathbb R^2\), where the fields have compact
support away from \(y_2=0\).  Because \(L\Psi\) is a divergence,
\(\int L\Psi=0\), and differentiation in time gives
\(\int\partial_tL\Psi=0\).  Also
\(v\cdot\nabla_y\xi=\operatorname{div}_y(v\xi)\), because
\(\operatorname{div}_yv=0\), so its integral is zero.  The function
\(W\) depends only on \(y_2\), and therefore
\(\int W\partial_{y_1}\Gamma^2=0\).  Finally,
\[
 \int M\xi\,dy
 =\int\bigl(\xi_{y_1y_1}+2y_2\xi_{y_2y_2}+4\xi_{y_2}\bigr)dy=0:
\]
the first and last terms integrate to zero, while integration by parts
in the middle term gives \(-2\int\xi_{y_2}=0\).
Thus \(\int E_\nu^\xi=0\), for Euler and for classical viscosity.
Lemma~\ref{lem:ec-inverse-div} gives
\(\operatorname{curl}_yb=E_\nu^\xi\), so
\eqref{eq:ec-force-recovery} satisfies both equations in
\eqref{eq:ec-reduced-NS}.  Its smoothness and support follow from that
lemma and from the fixed positive lower bound on \(r\) on
\(\widetilde D\).

It remains to construct the pressure explicitly rather than leave it
as a compatibility assertion.  Define the smooth physical vector field
\begin{equation}
 h=f+\nu\Delta_xu-\partial_tu-(u\cdot\nabla_x)u.
 \label{eq:ec-h}
\end{equation}
It is axisymmetric, compactly supported in \(\widetilde D\), and
identically zero in a fixed collar of the axis.  The circulation
equation gives \(h^\phi=0\).  The previously proved meridional
calculation gives the exact identity
\begin{equation}
 \partial_z h^r-\partial_r h^z
 =r\left[\operatorname{curl}_y(f^z,f^r/r)
        +\nu M\xi-D_v\xi+W\partial_{y_1}\Gamma^2\right]=0.
 \label{eq:ec-h-curl}
\end{equation}
This also specifies the precise compatibility map from the momentum
residual to the scalar residual, including the factor \(r\).

For \(r\geq0\), extend the components of \(h\) by zero at the axis
and put
\begin{equation}
 p(z,r,t)=\int_0^r h^r(z,\rho,t)\,d\rho.
 \label{eq:ec-pressure}
\end{equation}
The integrand vanishes in a fixed neighborhood of \(\rho=0\), so
there is no singular endpoint.  Direct differentiation yields
\[
 p_r=h^r,\qquad
 p_z=\int_0^r\partial_z h^r(z,\rho,t)\,d\rho
     =\int_0^r\partial_\rho h^z(z,\rho,t)\,d\rho=h^z(z,r,t).
\]
The last equality uses \(h^z(z,0,t)=0\).  There is no azimuthal
derivative, so \(\nabla_xp=h\) for \(r>0\).  Moreover \(p=0\)
in a fixed neighborhood of the axis.  It therefore extends smoothly
as an axisymmetric scalar on \(\mathbb R^3\), with all time
derivatives included.  Substitution of \eqref{eq:ec-h} proves
\eqref{eq:ec-NS} everywhere, including the axis.

For the converse, Proposition~\ref{prop:ec-velocity} recovers the
unique compact streamfunction and circulation from the original
velocity.  Applying the explicit component and curl calculations
already proved gives \eqref{eq:ec-reduced-NS}.  The original pressure
is allowed: no step replaced it by a force or changed its derivatives.
This proves the asserted correspondence.
\end{proof}

There are uniform finite derivative bounds in this reconstruction.
For nonnegative integers \(k,b\), the spatial chain and product rules,
Lemma~\ref{lem:ec-inverse-div}, and the fixed support give
\begin{equation}
 \|\partial_t^b f\|_{C_x^k}
 \leq C_k\left(
       \|\partial_t^b E_\nu^\xi\|_{C_y^k}
      +\|\partial_t^b F_\nu^\Gamma\|_{C_y^k}\right),
 \label{eq:ec-physical-force-bound}
\end{equation}
where \(C_k\) depends on the fixed geometry and bump, not on the
input pair or time.  To see explicitly why the coordinate factors are
bounded, on the compact support \(r\) has a positive minimum and a
finite maximum, so every finite derivative of \(r\), \(1/r\),
\(\Phi\), \(\Phi^{-1}\), \(e_r\), and \(e_\phi\) is bounded there.
Each derivative in the product and chain rules is a finite sum of
products of these fixed derivatives with derivatives of the indicated
input.  This proves \eqref{eq:ec-physical-force-bound} at every order.
The pressure formula likewise preserves smoothness at each order by
differentiation of its integral over a bounded radial interval on any
compact spatial set.

Force recovery is unique up to an exactly characterized gradient.
Suppose two smooth axisymmetric force fields supported away from the
axis realize the same \(F_\nu^\Gamma,E_\nu^\xi\).  Their difference
\(g\) has \(g^\phi=0\) and
\(\partial_zg^r-\partial_rg^z=0\).  The same formula
\(q(z,r,t)=\int_0^r g^r(z,\rho,t)\,d\rho\) then proves
\(g=\nabla_xq\).  Conversely adding \(\nabla_xq\) to the force and
\(q\) to the pressure preserves \eqref{eq:ec-NS} and both reduced
residuals, by direct substitution and cancellation of mixed
derivatives.  Thus the missing scalar in a curl formulation is an
explicit pressure gauge, with a proved reconstruction map.

\subsection{What the viscosity calculation supplies for the next extension}

At the level of arbitrary smooth fields, rather than just solutions,
the same pair \((\Gamma,\Psi)\) has residual differences
\begin{equation}
 F_\nu^\Gamma-F_0^\Gamma=-\nu Q\Gamma,\qquad
 E_\nu^\xi-E_0^\xi=-\nu M\xi.
 \label{eq:ec-viscous-difference}
\end{equation}
These identities exhibit the precise two additional terms that must
be controlled in the circulation and reduced-vorticity equations.
The coefficient \(2y_2=r^2\), the drift \(4\partial_{y_2}\), the
source coefficient \(W=r^{-4}\), and the original \(r_0,z_0,R,\nu\)
have all been retained.  The square equation
\eqref{eq:ec-Theta} also retains its nonzero viscous gradient-square
term.

There is an equally exact physical force identity.  For the same
smooth \((u,p)\), define
\(f_0=\partial_tu+(u\cdot\nabla)u+\nabla p\).  Then the unique
force that makes this same \((u,p)\) satisfy the equation with viscosity
\(\nu\) is
\begin{equation}
 f_\nu=f_0-\nu\Delta_xu.
 \label{eq:ec-physical-force-difference}
\end{equation}
This follows by rearranging \eqref{eq:ec-NS}; it fixes the sign.
The compact recovered force in \eqref{eq:ec-force-recovery} may use a
different pressure gauge, but its difference from this expression has
zero azimuthal component and zero meridional curl, by
\eqref{eq:ec-viscous-difference}.  The preceding explicit potential
construction proves that the difference is a smooth gradient.

The coordinate mechanism therefore supplies a concrete translation
between the equations.  It does not estimate
\(\partial_t^b\partial_y^\alpha Q\Gamma\),
\(\partial_t^b\partial_y^\alpha M\xi\), or
\(\partial_t^b\partial_x^\alpha\Delta_xu\) at the terminal time of
the released Euler construction.  Smoothness on each closed interval
strictly before that time gives no uniform terminal bound by itself:
even the scalar function \((T-t)^{-1}\) is smooth on every such
interval and is unbounded as \(t\uparrow T\).  A classical
Navier--Stokes construction requires actual estimates or cancellations
for the displayed viscous terms in its own complete evolution.
Those estimates have not been established in this section.

\subsection{Independent symbolic verification and its exact scope}

The accompanying script
\texttt{scripts/check\_euler\_coordinates.py} verifies fifteen exact
differential identities with arbitrary smooth symbolic functions
\(\Psi,\Gamma,p,f^r,f^\phi,f^z\), retaining \(r>0\) and
\(\nu\geq0\).  All fifteen remainders are the exact integer zero;
the recorded output is
\path{checks/euler_coordinates_symbolic.json}.
In particular, define the complete physical momentum residuals
\begin{align*}
 R_r&=Du^r-(u^\phi)^2/r+p_r
            -\nu(\Delta_0-r^{-2})u^r-f^r,\\
 R_\phi&=Du^\phi+u^ru^\phi/r
            -\nu(\Delta_0-r^{-2})u^\phi-f^\phi,\\
 R_z&=Du^z+p_z-\nu\Delta_0u^z-f^z.
\end{align*}
The independent expansion checks the identities
\begin{align}
 rR_\phi&=D\Gamma-\nu\mathcal D_-\Gamma-rf^\phi,
 \label{eq:ec-symbolic-swirl}\\
 \partial_zR_r-\partial_rR_z
 &=r\left[D\xi-r^{-4}\partial_z\Gamma^2-\nu\mathcal D_+\xi
               -\frac{\partial_zf^r-\partial_rf^z}{r}\right].
 \label{eq:ec-symbolic-curl}
\end{align}
Neither identity assumes that the residuals vanish.  The remaining
checks include the streamfunction signs, divergence, force-free
pressure cancellation, vector-Laplacian curl, both diffusion
coordinate changes, the material derivative, the vorticity operator,
the swirl source, and the oriented Jacobian.  These algebraic checks
corroborate the local differential calculation.  The global support,
smooth axis extension, inverse streamfunction, compact force inverse,
and pressure recovery are established by the proofs above, not by the
symbolic program.
